\section{Latihan}
\label{sec:latihan-bab2}

\begin{enumerate}
    \item \textbf{Konversi dasar.} Konversi bilangan berikut tanpa
          alat bantu: (a) \dec{173} ke biner dan heksadesimal;
          (b) \hex{5C} ke biner dan desimal; (c) \bin{1101\,0011} ke
          heksadesimal dan desimal.

    \item \textbf{\eng{Two's complement}.} Nyatakan bilangan berikut
          dalam 8-bit \eng{two's complement}: (a) \dec{-1};
          (b) \dec{-16}; (c) \dec{-100}. Tunjukkan langkah pembalikan
          bit dan penambahan satu untuk setiap soal.

    \item \textbf{Rentang nilai.} Berapa nilai terbesar dan terkecil
          yang dapat dinyatakan dalam 8 bit secara bertanda? Bagaimana
          bila dinyatakan secara tak bertanda? Jelaskan mengapa
          rentangnya berbeda.

    \item \textbf{Analisis flag.} Diberikan \reg{AL} = \hex{7F} dan
          instruksi \asm{ADD AL, 01h}. Tentukan hasil di \reg{AL}
          serta nilai \flag{CF}, \flag{ZF}, \flag{SF}, dan
          \flag{OF}. Jelaskan mengapa \flag{OF} bernilai 1 sementara
          \flag{CF} bernilai 0.

    \item \textbf{Analisis flag.} Ulangi soal 4 dengan \reg{AL} =
          \hex{FF}. Tentukan hasil dan seluruh flag, lalu jelaskan
          mengapa kali ini \flag{CF} bernilai 1 tetapi \flag{OF}
          bernilai 0.

    \item \textbf{Penyandian.} Sebuah variabel dideklarasikan sebagai
          \code{DB '5'}. Berapa nilai byte yang tersimpan? Berapa
          nilainya bila dideklarasikan sebagai \code{DB 5}? Jelaskan
          perbedaan dan akibatnya bila variabel itu dicetak dengan
          layanan DOS fungsi \hex{02}.

    \item \textbf{\eng{Little-endian}.} Sebuah variabel
          dideklarasikan \code{DW 0ABCDh}. Tuliskan urutan byte yang
          tersimpan di memori dari alamat terendah ke tertinggi, lalu
          jelaskan mengapa urutannya demikian.

    \item \textbf{Kaitan ke bab berikutnya.} Berdasarkan
          Gambar~\ref{fig:peta-konsep}, jelaskan mengapa
          Bab~\ref{chap:segmentasi} harus dipelajari sebelum
          Bab~\ref{chap:mode-pengalamatan}. Kaitkan jawaban Anda
          dengan konsep alamat efektif.
\end{enumerate}

\begin{refleksibox}
Jawab pertanyaan berikut pada jurnal praktikum Anda.
\begin{enumerate}
    \item Bagian mana dari representasi data yang paling menantang
          bagi Anda: konversi bilangan, \eng{two's complement}, atau
          pembedaan \eng{carry} dan \eng{overflow}? Mengapa?
    \item Setelah melihat isi memori secara langsung, apakah
          pemahaman Anda tentang ``penyimpanan data'' berubah?
          Bagian mana yang paling mengejutkan?
    \item Konsep apa pada bab ini yang Anda rasa akan paling sering
          Anda pakai pada bab-bab berikutnya?
\end{enumerate}
\end{refleksibox}
